\documentclass[answers,12pt,addpoints]{exam}

\usepackage{import}

% Use your existing style file if available.
% Otherwise, use a local style.tex or a basic margin setup.
\IfFileExists{C:/Users/prana/OneDrive/Desktop/MathNotes/style.tex}{%
  \import{C:/Users/prana/OneDrive/Desktop/MathNotes/}{style.tex}%
}{%
  \IfFileExists{style.tex}{%
    \input{style.tex}%
  }{%
    \usepackage[margin=1in]{geometry}%
  }%
}

\usepackage{amsmath,amssymb}
\usepackage{booktabs}
\usepackage{listings}
\usepackage{xcolor}
\usepackage{graphicx}

% Header
\providecommand{\name}{Pranav Tikkawar}
\providecommand{\course}{Stats 596}
\providecommand{\assignment}{Homework 2}

\author{\name}
\title{\course\ -- \assignment}
\date{Fall 2026\\Due Thursday, October 8, 2026}

% Notation
\providecommand{\R}{\mathbb{R}}
\providecommand{\E}{\operatorname{E}}
\providecommand{\argmin}{\operatorname*{arg\,min}}
\providecommand{\trans}{^{\mathsf T}}
\providecommand{\normtwo}[1]{\left\lVert #1\right\rVert_2}
\providecommand{\ones}{\mathbf{1}}
\providecommand{\ind}{\mathbf{1}}
\providecommand{\Var}{\operatorname{Var}}
\providecommand{\tr}{\operatorname{tr}}
\providecommand{\MSE}{\operatorname{MSE}}
\providecommand{\sgn}{\operatorname{sgn}}

% Reduce distracting overflow without forcing all paragraphs flush.
\setlength{\emergencystretch}{3em}

% R listing style
\definecolor{codegreen}{rgb}{0,0.45,0.15}
\definecolor{codepurple}{rgb}{0.50,0,0.50}
\definecolor{codegray}{rgb}{0.40,0.40,0.40}
\definecolor{codeback}{rgb}{0.97,0.97,0.97}

\lstdefinestyle{Rstyle}{
  language=R,
  backgroundcolor=\color{codeback},
  basicstyle=\ttfamily\small,
  keywordstyle=\color{blue},
  commentstyle=\color{codegreen},
  stringstyle=\color{codepurple},
  numberstyle=\tiny\color{codegray},
  numbers=left,
  numbersep=6pt,
  breaklines=true,
  columns=fullflexible,
  showstringspaces=false,
  keepspaces=true,
  tabsize=2
}
\lstset{style=Rstyle}

% Import numerical results generated by homework2.R.
% Missing files are clearly marked instead of inventing results.
% Absolute location of the results generated by the R script.
% Change ONLY this line if your results are in a different folder.
% Keep the trailing slash.
\newcommand{\resultsdir}{C:/Users/prana/OneDrive/Desktop/MathNotes/Stats 596/hw2_results/}

% Import generated LaTeX tables and summaries.
\newcommand{\resultinput}[1]{%
  \IfFileExists{\resultsdir#1}{%
    \input{\resultsdir#1}%
  }{%
    \begin{center}
    \fbox{%
      \parbox{0.88\linewidth}{%
        \raggedright
        Result file not found.
        \par\smallskip
        Expected folder:
        \texttt{\expandafter\detokenize\expandafter{\resultsdir}}
        \par\smallskip
        Expected file:
        \texttt{\detokenize{#1}}
      }%
    }
    \end{center}
  }%
}

% Import generated PDF figures.
\newcommand{\resultfigure}[2]{%
  \begin{center}
  \IfFileExists{\resultsdir#1}{%
    \includegraphics[width=0.98\linewidth]{\resultsdir#1}%
  }{%
    \fbox{%
      \parbox{0.88\linewidth}{%
        \raggedright
        Figure file not found.
        \par\smallskip
        Expected folder:
        \texttt{\expandafter\detokenize\expandafter{\resultsdir}}
        \par\smallskip
        Expected file:
        \texttt{\detokenize{#1}}
      }%
    }%
  }
  \par\smallskip
  {\small #2}
  \end{center}
}

\begin{document}
\maketitle

\begin{questions}

%==============================================================================
\question \textbf{Exercise 1.11.}

Let $Z\sim N(\mu,1)$, where $\mu\in\R$.

\begin{parts}

\part For $\lambda\geq0$, define the ridge estimator
\[
\widehat\mu_\lambda(z)
=
\argmin_{t\in\R}
\left\{
\frac12(t-z)^2+\frac{\lambda}{2}t^2
\right\}.
\]
Compute
\[
\MSE_\lambda(\mu)
=
\E\{(\widehat\mu_\lambda(Z)-\mu)^2\}.
\]
Plot and compare the risk functions for
$\lambda=0.1,1,5,10$.

\part Repeat for the Lasso estimator
\[
\widehat\mu_\lambda(z)
=
\argmin_{t\in\R}
\left\{
\frac12(t-z)^2+\lambda|t|
\right\}.
\]

\end{parts}

\begin{solution}

\textbf{(a) Ridge.}
The derivative of the objective is
\[
(t-z)+\lambda t=(1+\lambda)t-z.
\]
Its second derivative is $1+\lambda>0$, so the unique minimizer is
\[
\boxed{
\widehat\mu_\lambda(z)=\frac{z}{1+\lambda}.
}
\]

Since $\E Z=\mu$ and $\Var(Z)=1$,
\[
\E\widehat\mu_\lambda(Z)-\mu
=
-\frac{\lambda\mu}{1+\lambda},
\qquad
\Var\{\widehat\mu_\lambda(Z)\}
=
\frac{1}{(1+\lambda)^2}.
\]
Therefore
\[
\boxed{
\MSE_\lambda^{\mathrm{ridge}}(\mu)
=
\frac{1+\lambda^2\mu^2}{(1+\lambda)^2}.
}
\]

Each risk curve is even in $\mu$ and has its minimum at zero.
Larger penalties reduce variance near zero but increase shrinkage
bias for sufficiently large signals.

The unpenalized estimator $Z$ has MSE 1. For $\lambda>0$,
\[
\MSE_\lambda^{\mathrm{ridge}}(\mu)<1
\iff
1+\lambda^2\mu^2<(1+\lambda)^2
\iff
\boxed{
\mu^2<1+\frac{2}{\lambda}.
}
\]
Thus a positive ridge penalty improves on $Z$ within a bounded
range of signal sizes.

At $\mu=0$, the risks for the four penalties are approximately
\[
0.826446,\qquad
0.25,\qquad
0.027778,\qquad
0.008264.
\]

\textbf{(b) Lasso.}
The stationary condition is
\[
0\in t-z+\lambda\partial|t|.
\]
Considering positive, zero, and negative $t$ gives
\[
\boxed{
\widehat\mu_\lambda(z)
=
S_\lambda(z)
=
\begin{cases}
z-\lambda, & z>\lambda,\\
0, & |z|\leq\lambda,\\
z+\lambda, & z<-\lambda.
\end{cases}
}
\]

Write $Z=\mu+U$, where $U\sim N(0,1)$, and define
\[
a=-\lambda-\mu,
\qquad
b=\lambda-\mu.
\]
Then
\[
\widehat\mu_\lambda(Z)-\mu
=
\begin{cases}
U+\lambda, & U<a,\\
-\mu, & a\leq U\leq b,\\
U-\lambda, & U>b.
\end{cases}
\]

Let $\phi$ and $\Phi$ denote the standard normal density and CDF.
The MSE is
\begin{align*}
\MSE_\lambda^{\mathrm{lasso}}(\mu)
={}&
\int_{-\infty}^{a}(u+\lambda)^2\phi(u)\,du\\
&+\mu^2\{\Phi(b)-\Phi(a)\}\\
&+\int_b^\infty(u-\lambda)^2\phi(u)\,du.
\end{align*}

Using $\phi'(u)=-u\phi(u)$ and integration by parts,
\begin{align*}
\int_{-\infty}^{a}u\phi(u)\,du
&=-\phi(a),\\
\int_{-\infty}^{a}u^2\phi(u)\,du
&=\Phi(a)-a\phi(a),\\
\int_b^\infty u\phi(u)\,du
&=\phi(b),\\
\int_b^\infty u^2\phi(u)\,du
&=1-\Phi(b)+b\phi(b).
\end{align*}

Expanding the squares and substituting yields
\[
\boxed{
\begin{aligned}
\MSE_\lambda^{\mathrm{lasso}}(\mu)
={}&
(1+\lambda^2)\{\Phi(a)+1-\Phi(b)\}\\
&-(a+2\lambda)\phi(a)\\
&+(b-2\lambda)\phi(b)\\
&+\mu^2\{\Phi(b)-\Phi(a)\}.
\end{aligned}
}
\]

At $\lambda=0$, the formula reduces to 1.
At $\mu=0$, it becomes
\[
\MSE_\lambda^{\mathrm{lasso}}(0)
=
2\{(1+\lambda^2)\Phi(-\lambda)-\lambda\phi(\lambda)\}.
\]

The risk is even in $\mu$. Large penalties produce small risk near
zero by setting the estimator to zero with high probability.
For fixed $\lambda$,
\[
\boxed{
\lim_{|\mu|\to\infty}
\MSE_\lambda^{\mathrm{lasso}}(\mu)
=
1+\lambda^2.
}
\]
The limiting risks for the four penalties are
\[
1.01,\qquad 2,\qquad 26,\qquad 101.
\]
Unlike ridge, the Lasso risk approaches a finite plateau for each
fixed penalty.

\resultfigure{ex111_mse.pdf}{
Exact ridge and Lasso MSE curves for
$\lambda=0.1,1,5,10$.
The dashed line is the unpenalized risk 1.
}

The R script evaluates these formulas and independently checks
the Lasso expression by numerical integration.

\end{solution}

%==============================================================================
\question \textbf{Exercise 1.12.}

Let $Y_i$ be responses and let $X_i=X_{i,1:p}$ denote covariates
excluding a constant. Define
\[
\widetilde Y_i=Y_i-\overline Y,
\qquad
\widetilde X_i=X_i-\overline X.
\]

\begin{parts}

\part For ridge regression with an unpenalized intercept, show that
the slopes can be obtained from the centered data and that
\[
\widehat\beta_0
=
\overline Y-\overline X\trans\widetilde\beta_{1:p}.
\]
Also show that
\[
\tr(S_\lambda)=1+\tr(\widetilde S_\lambda).
\]

\part Prove the same centering and intercept identities for Lasso.

\end{parts}

\begin{solution}

\textbf{(a)}
Write $b=\beta_{1:p}$.
The ridge criterion is
\[
Q(\beta_0,b)
=
\frac{1}{2n}
\sum_{i=1}^n
(Y_i-\beta_0-X_i\trans b)^2
+
\frac{\lambda}{2}\normtwo{b}^2.
\]

For fixed $b$, the penalty is independent of the intercept.
Differentiating in $\beta_0$ gives
\[
\frac{\partial Q}{\partial\beta_0}
=
-\frac1n
\sum_{i=1}^n
(Y_i-\beta_0-X_i\trans b).
\]
The second derivative is 1, so the unique minimizing intercept is
\[
\beta_0(b)=\overline Y-\overline X\trans b.
\]

Substituting into the residual gives
\[
Y_i-\beta_0(b)-X_i\trans b
=
\widetilde Y_i-\widetilde X_i\trans b.
\]
Consequently the profiled objective is
\[
\min_{\beta_0}Q(\beta_0,b)
=
\frac{1}{2n}
\sum_{i=1}^n
(\widetilde Y_i-\widetilde X_i\trans b)^2
+
\frac{\lambda}{2}\normtwo{b}^2.
\]
This is exactly the centered ridge problem. Therefore
\[
\boxed{
\widehat\beta_{1:p}=\widetilde\beta_{1:p},
\qquad
\widehat\beta_0
=
\overline Y-\overline X\trans\widetilde\beta_{1:p}.
}
\]
Any centered minimizer gives a joint minimizer, even when the slope
solution is not unique at $\lambda=0$.

For the smoother identity, let $X_c$ be the original $n\times p$
covariate matrix and define
\[
J=\frac1n\ones\ones\trans,
\qquad
H=I-J.
\]
Then
\[
\widetilde X=HX_c,
\qquad
\widetilde Y=HY.
\]

For $\lambda>0$, the centered smoother is
\[
\widetilde S_\lambda
=
\widetilde X
(\widetilde X\trans\widetilde X+n\lambda I)^{-1}
\widetilde X\trans.
\]
Because $\widetilde X\trans\ones=0$,
\[
\widetilde S_\lambda J=0,
\qquad
\widetilde S_\lambda H=\widetilde S_\lambda.
\]

The fitted response from the original regression is
\begin{align*}
\widehat Y
&=
\ones\widehat\beta_0+X_c\widetilde\beta_{1:p}\\
&=
\overline Y\ones+\widetilde X\widetilde\beta_{1:p}\\
&=
JY+\widetilde S_\lambda HY\\
&=
(J+\widetilde S_\lambda)Y.
\end{align*}
Thus
\[
S_\lambda=J+\widetilde S_\lambda.
\]
Since $\tr(J)=1$,
\[
\boxed{
\tr(S_\lambda)=1+\tr(\widetilde S_\lambda).
}
\]

At $\lambda=0$, use the centered least-squares projection matrix,
defined with a Moore--Penrose inverse when necessary.
The same fitted-value argument applies.

\textbf{(b)}
For Lasso, the objective becomes
\[
Q(\beta_0,b)
=
\frac{1}{2n}
\sum_{i=1}^n
(Y_i-\beta_0-X_i\trans b)^2
+
\lambda\|b\|_1.
\]
The penalty still does not involve $\beta_0$, so the optimal
intercept and residual substitution are unchanged. Hence
\[
\boxed{
\widehat\beta_{1:p}=\widetilde\beta_{1:p},
\qquad
\widehat\beta_0
=
\overline Y-\overline X\trans\widetilde\beta_{1:p}.
}
\]
No differentiation of the Lasso penalty is needed.

\end{solution}

%==============================================================================
\question \textbf{Exercise 1.13.}

Let
\[
s_j=\|X_j\|_n
=
\left(\frac1n\sum_{i=1}^nX_{ij}^2\right)^{1/2},
\qquad
\widetilde X_{ij}=\frac{X_{ij}}{s_j}.
\]

\begin{parts}

\part Find weights $c_j$ so that fitting an unweighted ridge penalty
to the standardized covariates and setting
$\widehat\beta_j=\widehat\gamma_j/s_j$ is equivalent to fitting
a weighted ridge penalty to the original covariates.

\part Find the corresponding weights for Lasso.

\end{parts}

\begin{solution}

Assume $s_j>0$ for each $j$, since otherwise the stated
standardization is undefined.

Let
\[
D=\operatorname{diag}(s_1,\ldots,s_p).
\]
Then $\widetilde X=XD^{-1}$, and the change of variables
\[
\gamma=D\beta
\]
is a bijection satisfying
\[
\widetilde X\gamma=X\beta.
\]

\textbf{(a) Ridge.}
The standardized objective becomes
\begin{align*}
&
\frac1{2n}\|Y-\widetilde X\gamma\|_2^2
+\frac{\lambda}{2}\sum_{j=1}^p\gamma_j^2\\
&\qquad=
\frac1{2n}\|Y-X\beta\|_2^2
+\frac{\lambda}{2}\sum_{j=1}^p s_j^2\beta_j^2.
\end{align*}
The objectives agree under the bijection, so their minimizers
correspond. Therefore
\[
\boxed{
c_j=s_j^2=\frac1n\sum_{i=1}^nX_{ij}^2.
}
\]

\textbf{(b) Lasso.}
Similarly,
\begin{align*}
&
\frac1{2n}\|Y-\widetilde X\gamma\|_2^2
+\lambda\sum_{j=1}^p|\gamma_j|\\
&\qquad=
\frac1{2n}\|Y-X\beta\|_2^2
+\lambda\sum_{j=1}^p s_j|\beta_j|.
\end{align*}
Thus
\[
\boxed{
c_j=s_j
=
\left(\frac1n\sum_{i=1}^nX_{ij}^2\right)^{1/2}.
}
\]

Standardization therefore changes the relative penalty weights
on the original coefficients.
At $\lambda=0$, the penalty weights are immaterial, but the stated
choices remain valid.

\end{solution}

%==============================================================================
\question \textbf{Exercise 1.15.}

\begin{parts}

\part For the Lasso objective
\[
h(\beta)
=
\frac1{2n}\|Y-X\beta\|_2^2+\lambda\|\beta\|_1,
\]
show that coordinate-wise minimization implies joint minimization.

\part Generalize to
\[
h(\beta)=h_0(\beta)+\sum_{j=1}^p h_{1,j}(\beta_j),
\]
where $h_0$ is convex and everywhere differentiable, and each
$h_{1,j}$ is convex but may be nondifferentiable.

\part (Optional.) Give a convex counterexample for which
coordinate-wise minimization does not imply joint minimization.

\end{parts}

\begin{solution}

\textbf{(a)}
Let
\[
h_0(\beta)=\frac1{2n}\|Y-X\beta\|_2^2.
\]
Its gradient is
\[
\nabla h_0(\beta)
=
-\frac1nX\trans(Y-X\beta).
\]

Suppose $\widetilde\beta$ minimizes the objective in every
coordinate when the other coordinates are held fixed.
The one-dimensional optimality condition for coordinate $j$ is
\[
0\in
\frac{\partial h_0}{\partial\beta_j}(\widetilde\beta)
+\lambda\partial|\widetilde\beta_j|.
\]

The absolute-value subdifferential is
\[
\partial|t|
=
\begin{cases}
\{1\}, & t>0,\\
{[-1,1]}, & t=0,\\
\{-1\}, & t<0.
\end{cases}
\]

For each $j$, select
$s_j\in\partial|\widetilde\beta_j|$ such that
\[
\frac{\partial h_0}{\partial\beta_j}(\widetilde\beta)
+\lambda s_j=0.
\]
Since the $L_1$ penalty is separable,
\[
s=(s_1,\ldots,s_p)\trans
\in\partial\|\widetilde\beta\|_1.
\]
Combining the coordinate conditions gives
\[
0
=
\nabla h_0(\widetilde\beta)+\lambda s
\in\partial h(\widetilde\beta).
\]

By the subgradient inequality, for every $\beta$,
\[
h(\beta)
\geq
h(\widetilde\beta)
+
0\trans(\beta-\widetilde\beta)
=
h(\widetilde\beta).
\]
Therefore $\widetilde\beta$ is a global minimizer.

Equivalently, its stationary conditions are
\[
\frac1nX_j\trans(Y-X\widetilde\beta)
\begin{cases}
=\lambda, & \widetilde\beta_j>0,\\
\in[-\lambda,\lambda], & \widetilde\beta_j=0,\\
=-\lambda, & \widetilde\beta_j<0.
\end{cases}
\]

This establishes the optimality of a coordinate-wise minimizing
point; it is not by itself a proof of convergence of every
coordinate-descent implementation.

\textbf{(b)}
Coordinate-wise optimality implies that, for each $j$, there exists
\[
s_j\in\partial h_{1,j}(\widetilde\beta_j)
\]
satisfying
\[
s_j=-\frac{\partial h_0}{\partial\beta_j}(\widetilde\beta).
\]

Convexity of $h_0$ gives
\[
h_0(\beta)
\geq
h_0(\widetilde\beta)
+
\nabla h_0(\widetilde\beta)\trans
(\beta-\widetilde\beta).
\]
The subgradient inequality for each penalty term gives
\[
h_{1,j}(\beta_j)
\geq
h_{1,j}(\widetilde\beta_j)
+
s_j(\beta_j-\widetilde\beta_j).
\]

Adding these inequalities,
\begin{align*}
h(\beta)
&\geq
h(\widetilde\beta)
+
\sum_{j=1}^p
\left\{
\frac{\partial h_0}{\partial\beta_j}(\widetilde\beta)
+s_j
\right\}
(\beta_j-\widetilde\beta_j)\\
&=
h(\widetilde\beta).
\end{align*}
Hence $\widetilde\beta$ is jointly minimizing.
Separability allows the individual penalty subgradients to
be chosen simultaneously.

\textbf{(c)}
Consider
\[
\boxed{
h(x,y)=|x-y|+\frac14(x+y-2)^2.
}
\]
The function is convex because it is the sum of the absolute
value of a linear function and the square of an affine function.

At $(0,0)$, fixing $y=0$ gives
\[
h(x,0)-h(0,0)
=
|x|-x+\frac{x^2}{4}
\geq0.
\]
Equality holds only at $x=0$.
By symmetry,
\[
h(0,y)-h(0,0)
=
|y|-y+\frac{y^2}{4}
\geq0,
\]
with equality only at $y=0$.

Thus $(0,0)$ uniquely minimizes the objective along each
coordinate line. However,
\[
h(0,0)=1,
\qquad
h(1,1)=0.
\]
Since $h\geq0$, $(1,1)$ is a global minimizer and $(0,0)$ is not.

The nondifferentiable term $|x-y|$ couples the coordinates;
it does not satisfy the separability assumption in part (b).

\end{solution}

%==============================================================================
\question \textbf{Exercise 1.17.}

Generate training and test data using the following code.

\begin{lstlisting}
library(MASS)
set.seed(0)

n <- 200
p <- 20
mu <- rep(0,p)
msig <- toeplitz(.8^(1:p-1))
sig <- .5

x <- mvrnorm(n, mu, msig)
y <- x[,1] +
  pmax(0, apply(x[,2:4],1,sum))^2 / 3 +
  rnorm(n, 0, sig)

cbind(y,x)[1:3, 1:6]

ntest <- 100
xtest <- mvrnorm(ntest, mu, msig)
ytest <- xtest[,1] +
  pmax(0, apply(xtest[,2:4],1,sum))^2 / 3 +
  rnorm(ntest, 0, sig)

cbind(ytest,xtest)[1:3, 1:6]
\end{lstlisting}

Apply all methods without standardizing $x$ or $y$.

\begin{parts}

\part Fit ridge regression with an unpenalized intercept and
produce a coefficient path.

\part Plot training, test, 5-fold CV, and 10-fold CV errors
against degrees of freedom. Report the minimum test error and
test errors at the minimum-CV and one-standard-error choices.

\part Fit Lasso with an unpenalized intercept and produce
a coefficient path.

\part Plot the four errors against the shrinkage fraction and
report the corresponding test-error comparisons.

\part Summarize the findings in one paragraph.

\end{parts}

\begin{solution}

\textbf{Data and reproducibility.}
The covariance is
\[
\Sigma_{jk}=0.8^{|j-k|},
\]
and the generating conditional mean is
\[
m(X)
=
X_1+\frac13\{\max(0,X_2+X_3+X_4)\}^2.
\]

The first three training rows, showing the response and the
first five covariates, are
\resultinput{p20_train_rows.tex}

The corresponding test rows are
\resultinput{p20_test_rows.tex}

The R script generates both samples before generating CV folds.
It then uses a separate seed of 123 for the fold partition.
Ridge and Lasso use the same partitions.
Neither $x$ nor $y$ is rescaled; centering profiles out
the unpenalized intercept.

\textbf{(a) Ridge.}
For each penalty,
\[
\widehat b_\lambda
=
(\widetilde X\trans\widetilde X+n\lambda I)^{-1}
\widetilde X\trans\widetilde Y,
\]
and
\[
\widehat\beta_{0,\lambda}
=
\overline Y-\overline X\trans\widehat b_\lambda.
\]

The implementation uses an SVD rather than explicitly
inverting the Gram matrix.
If $d_1,\ldots,d_r$ are the retained nonzero singular values
of $\widetilde X$, the effective degrees of freedom are
\[
\boxed{
\operatorname{df}(\lambda)
=
1+\sum_{j=1}^r
\frac{d_j^2}{d_j^2+n\lambda}.
}
\]
The intercept contributes the leading 1.

The grid contains 240 positive penalties from
$10^4s$ to $10^{-8}s$, where
\[
s=\frac{d_{\max}^2}{n}.
\]
This sets the penalty-grid scale without standardizing the data.

\resultfigure{p20_ridge_path.pdf}{
Ridge coefficient path versus effective degrees of freedom.
}

\textbf{(b) Errors and cross-validation.}
Training and test errors are mean squared prediction errors:
\[
\operatorname{Err}_{\rm train}(\lambda)
=
\frac1n\sum_{i=1}^n
(Y_i-\widehat Y_{i,\lambda})^2,
\]
\[
\operatorname{Err}_{\rm test}(\lambda)
=
\frac1{n_{\rm test}}\sum_{i=1}^{n_{\rm test}}
(Y_i^{\rm test}-\widehat Y_{i,\lambda}^{\rm test})^2.
\]

For $K\in\{5,10\}$, let $e_k(\lambda)$ be the mean squared
held-out error in fold $k$, using a fit trained on the other folds.
The folds have equal size, so
\[
\operatorname{CV}_K(\lambda)
=
\frac1K\sum_{k=1}^K e_k(\lambda),
\]
and the fold-based standard error is
\[
\operatorname{SE}_K(\lambda)
=
\frac{
\operatorname{sd}\{e_1(\lambda),\ldots,e_K(\lambda)\}
}{\sqrt K}.
\]

The minimum-CV choice is
\[
\lambda_{\min,K}
\in
\argmin_{\lambda\in\Lambda}
\operatorname{CV}_K(\lambda).
\]
The one-standard-error choice is the largest penalty satisfying
\[
\operatorname{CV}_K(\lambda)
\leq
\operatorname{CV}_K(\lambda_{\min,K})
+
\operatorname{SE}_K(\lambda_{\min,K}).
\]
Thus the one-SE rule favors stronger regularization among
fits whose CV error is sufficiently close to the minimum.

This standard error is a fold-based selection heuristic;
overlapping training sets do not produce independent fits.

\resultfigure{p20_ridge_errors.pdf}{
Ridge training, test, and CV errors.
The plots show the full error range and a detailed view.
}

\resultfigure{p20_ridge_cv.pdf}{
Ridge CV means with one-SE bands and selected penalties.
}

\resultinput{p20_ridge_selection.tex}

The test-grid oracle is the smallest observed test error among
the 240 fitted penalties. It is reported only as a benchmark:
test data are not used to select the CV penalties.
A minimum over a finite grid need not equal the exact minimum
over all $\lambda\geq0$.

\textbf{(c) Lasso.}
The objective is
\[
\frac1{2n}
\|Y-\beta_0\ones-Xb\|_2^2+\lambda\|b\|_1.
\]

The implementation uses \texttt{glmnet} with these settings:
\begin{lstlisting}
alpha = 1
standardize = FALSE
intercept = TRUE
\end{lstlisting}

The script checks the original-scale KKT conditions, including
the zero-mean residual condition for the intercept.
For Gaussian Lasso, the package's internal response transformation
is undone and does not change this objective.
The script itself does not standardize the response.
Ridge is implemented separately to preserve its stated
original-scale tuning convention.

Define
\[
\lambda_{\max}
=
\max_j
\left|
\frac1n\widetilde X_j\trans\widetilde Y
\right|.
\]
At this penalty, zero slopes satisfy the stationary conditions.
The Lasso grid has 240 log-spaced penalties decreasing from
$\lambda_{\max}$ to $10^{-6}\lambda_{\max}$.

The shrinkage fraction is
\[
f(\lambda)
=
\frac{\|\widehat b_\lambda\|_1}
{\|\widehat b_{\rm LS}\|_1}.
\]
For $p=20$, the centered design has full column rank almost
surely, so the least-squares slope reference is unique.

\resultfigure{p20_lasso_path.pdf}{
Lasso coefficient path versus the $L_1$ shrinkage fraction.
}

\textbf{(d) Lasso errors.}
Use the same MSE definitions, folds, and selection rules as
in part (b), now plotting against $f(\lambda)$.

\resultfigure{p20_lasso_errors.pdf}{
Lasso training, test, and CV errors.
}

\resultfigure{p20_lasso_cv.pdf}{
Lasso CV means with one-SE bands and selected penalties.
}

\resultinput{p20_lasso_selection.tex}

The script exports every penalty, error value, and fitted
coefficient to CSV files for reproducibility.

\textbf{(e) Findings.}
\resultinput{ex117_summary.tex}

\end{solution}

%==============================================================================
\question \textbf{Exercise 1.18.}

Repeat Exercise 1.17 with $p=100$ and $p=200$.
For each case, answer questions (i)--(iv) of Exercise 1.17,
then summarize the findings in one paragraph.

\begin{solution}

Each dimension uses the same simulation code as Exercise 1.17,
with only $p$ changed. The seed is reset to zero before each
complete training/test simulation, and the CV seed is applied
afterward. The fitting methods, error definitions, and
selection rules are unchanged.

\textbf{The $p=100$ case.}

The first training and test rows are
\resultinput{p100_train_rows.tex}
\resultinput{p100_test_rows.tex}

The ridge results are
\resultfigure{p100_ridge_path.pdf}{
Ridge coefficient path, $p=100$.
}
\resultfigure{p100_ridge_errors.pdf}{
Ridge training, test, and CV errors, $p=100$.
}
\resultfigure{p100_ridge_cv.pdf}{
Ridge CV bands and selected penalties, $p=100$.
}
\resultinput{p100_ridge_selection.tex}

The Lasso results are
\resultfigure{p100_lasso_path.pdf}{
Lasso coefficient path, $p=100$.
}
\resultfigure{p100_lasso_errors.pdf}{
Lasso training, test, and CV errors, $p=100$.
}
\resultfigure{p100_lasso_cv.pdf}{
Lasso CV bands and selected penalties, $p=100$.
}
\resultinput{p100_lasso_selection.tex}

\textbf{The $p=200$ case.}

The first training and test rows are
\resultinput{p200_train_rows.tex}
\resultinput{p200_test_rows.tex}

The ridge results are
\resultfigure{p200_ridge_path.pdf}{
Ridge coefficient path, $p=200$.
}
\resultfigure{p200_ridge_errors.pdf}{
Ridge training, test, and CV errors, $p=200$.
}
\resultfigure{p200_ridge_cv.pdf}{
Ridge CV bands and selected penalties, $p=200$.
}
\resultinput{p200_ridge_selection.tex}

The Lasso results are
\resultfigure{p200_lasso_path.pdf}{
Lasso coefficient path, $p=200$.
}
\resultfigure{p200_lasso_errors.pdf}{
Lasso training, test, and CV errors, $p=200$.
}
\resultfigure{p200_lasso_cv.pdf}{
Lasso CV bands and selected penalties, $p=200$.
}
\resultinput{p200_lasso_selection.tex}

\textbf{Rank and reference convention.}
At $p=200$, centering $n=200$ observations gives
\[
\operatorname{rank}(\widetilde X)\leq199.
\]
The unpenalized least-squares slope solution is therefore
nonunique. The script uses the minimum-Euclidean-norm solution
\[
\widehat b_{\rm LS}
=
\widetilde X^+\widetilde Y
\]
as the explicit reference for the Lasso shrinkage fraction.

Different least-squares slope solutions can have different
$L_1$ norms, so this convention matters.
With this reference, a shrinkage fraction above 1 is possible;
it is not automatically a coding error.
The plots preserve penalty-path order rather than assuming
the fraction is monotone.

Positive ridge penalties remain well defined even within CV
folds whose training sample size is smaller than the number
of covariates.

Changing $p$ changes the realized observations, even when
the seed is the same. Thus these comparisons are not paired
comparisons on identical responses, although the generating
mean formula remains unchanged.

\textbf{Findings.}
\resultinput{ex118_summary.tex}

\end{solution}

\end{questions}

%==============================================================================
\clearpage
\section*{Reproducible R code}

Place \texttt{homework2.R} in the same folder as this document.
Run the following commands from the R console:

\begin{lstlisting}
setwd("C:/Users/prana/OneDrive/Desktop/MathNotes/Stats 596")

if (!requireNamespace("MASS", quietly = TRUE)) {
  install.packages("MASS")
}

if (!requireNamespace("glmnet", quietly = TRUE)) {
  install.packages("glmnet")
}

source("homework2.R")
\end{lstlisting}

The script creates the folder \texttt{hw2\_results}, writes
PDF figures, CSV tables, and LaTeX result fragments, and records
the R session information.
After it finishes, recompile this document.

Inspect any warnings about minima occurring at penalty-grid
endpoints before treating those minima as adequately resolved.
If numerical output has not yet been generated, this document
shows explicit placeholders instead of fabricated results.

\IfFileExists{homework2.R}{%
  \lstinputlisting{homework2.R}%
}{%
  \IfFileExists{hw2.R}{%
    \lstinputlisting{hw2.R}%
  }{%
    \begin{center}
    \fbox{%
      \parbox{0.88\linewidth}{%
        \raggedright
        R script not found.
        Place \texttt{homework2.R} beside this document
        to include the complete code listing.
      }%
    }
    \end{center}
  }%
}

\end{document}